% MT063 / explanation — canonical editable content.
先按残差\(r=y-f(x)\)画损失，不必先想权重。该函数在\(|r|<1\)时使用半平方误差，两端改为绝对值减0.5。


\textbf{(a) 先算关键点再连线。} \(L(0)=0\)；到\(r=\pm1\)时，中心公式给\(1/2\)，外侧公式也给\(1-1/2=1/2\)，所以两段接得上。到\(r=\pm2\)时，\(L=2-1/2=1.5\)。中心画开口向上的抛物线，\(-1\)左侧及1右侧画直线，不要继续画成抛物线。


\StudyImage{CS5489/exam-review/Midterm/assets/answer-huber.png}


它们的斜率为



\[
L'(r)=\begin{cases}-1,&r<-1,\\r,&-1\le r\le1,\\1,&r>1.\end{cases}
\]


中心到接点的斜率恰为\(\pm1\)，与两端直线相同，因此连接处没有折角。0附近也平滑，比纯绝对损失在0处的尖点更容易处理。


\textbf{(b) 为什么对大输出异常值更稳健。} 对残差10，半平方损失为50，斜率为10；Huber损失为9.5，斜率只有1。梯度更新根据斜率调整参数，Huber不会让这个极端误差仅因更远就无限增大拉动强度。


说得更具体：因\(r=y-w^Tx-b\)，对\(w\)求导得到\(-L'(r)x\)。Huber限制的是残差方向的\(L'(r)\)，输入\(x\)若极大仍会放大参数梯度，因此它主要缓和大输出残差，而不是保证对所有异常点都无影响。


\textbf{原题差异 / Source note：} 随卷答案图高度为题干公式的两倍；本图按公式，L(\ensuremath{\pm}1)=0.5。 / The supplied plot is scaled by two; this plot follows the formula.


\textbf{出处：} 2023B期中 Q12，原卷第9页（10分）


\textbf{答案依据：} 2023B期中 Q12 随卷／配套参考答案；条件与解释经复核。
